\clearpage

\section{Порядок работы системы}

Вопрос: \texttt{?- factorial(2, Res).}

\begin{tabularx}{\textwidth}{|c|>{\small}X|>{\small}X|>{\small}X|}
	\hline
	\textbf{№}                                                                         &
	\textbf{Состояние резольвенты}                                                     &
	\textbf{Для каких термов запускается алгоритм унификации, результат и подстановка} &
	\textbf{Дальнейшие действия}                                                               \\
	\hline

	0                                                                                  &
	factorial(2, Res).                                                                 &
	—                                                                                  &
	Резольвента непустая. \newline
	Прямой ход. \newline
	Запуск алгоритма унификации для верхней подцели с 1-ым заголовком из БЗ (factorial\_tail). \\
	\hline

	1                                                                                  &
	factorial(2, Res).                                                                 &
	factorial(2, Res) = \newline
	factorial\_tail(0, Acc, Acc) \newline
	Унификация неуспешна (разные функторы).                                            &
	Унификация неуспешна. \newline
	Резольвента непустая. \newline
	Прямой ход.                                                                                \\
	\hline

	2                                                                                  &
	...                                                                                &
	... (Пропуск второго правила factorial\_tail)                                      &
	...                                                                                        \\
	\hline

	3                                                                                  &
	factorial\_tail(2, 1, Res). \newline
	Успешная унификация. \newline
	Редукция и применение подстановки к полученной конъюнкции целей: \newline
	factorial\_tail(2, 1, Res).                                                        &
	factorial(2, Res) = \newline
	factorial(N, Result) \newline
	Унификация успешна. \newline
	Подстановка: \newline
	\{N = 2, Result = Res\}                                                            &
	Унификация успешна. \newline
	Резольвента непустая. \newline
	Прямой ход. \newline
	Запуск алгоритма унификации для factorial\_tail(2, 1, Res) и 1-ого заголовка.              \\
	\hline

	4                                                                                  &
	factorial\_tail(2, 1, Res).                                                        &
	factorial\_tail(2, 1, Res) = \newline
	factorial\_tail(0, Acc, Acc) \newline
	Унификация неуспешна (разные константы: 2 и 0).                                    &
	Унификация неуспешна. \newline
	Резольвента непустая. \newline
	Прямой ход. Переход к следующему правилу.                                                  \\
	\hline

	5                                                                                  &
	>(2, 0), \newline
	NewAcc is 1 * 2, \newline
	NewN is 2 - 1, \newline
	factorial\_tail(NewN, NewAcc, Res). \newline
	Успешная унификация. \newline
	Редукция и применение подстановки.                                                 &
	factorial\_tail(2, 1, Res) = \newline
	factorial\_tail(N, Acc, Result) \newline
	Унификация успешна. \newline
	Подстановка: \newline
	\{N = 2, Acc = 1, Result = Res\}                                                   &
	Унификация успешна. \newline
	Резольвента непустая. \newline
	Прямой ход. \newline
	Запуск алгоритма унификации для >(2, 0) и Истина.                                          \\
	\hline

	6                                                                                  &
	NewAcc is 1 * 2, \newline
	NewN is 2 - 1, \newline
	factorial\_tail(NewN, NewAcc, Res). \newline
	Успешная унификация. \newline
	Вычисление предиката сравнения.                                                    &
	>(2, 0) = \newline
	Истина \newline
	Унификация успешна. \newline
	Подстановка: \newline
	Нет.                                                                               &
	Унификация успешна. \newline
	Резольвента непустая. \newline
	Прямой ход. \newline
	Запуск алгоритма вычисления для is(1 * 2).                                                 \\
	\hline

	7                                                                                  &
	NewN is 2 - 1, \newline
	factorial\_tail(NewN, 2, Res). \newline
	Успешная унификация. \newline
	Вычисление предиката is.                                                           &
	NewAcc is 1 * 2 = \newline
	Истина \newline
	Унификация успешна. \newline
	Подстановка: \newline
	\{NewAcc = 2\}                                                                     &
	Унификация успешна. \newline
	Резольвента непустая. \newline
	Прямой ход. \newline
	Запуск алгоритма вычисления для is(2 - 1).                                                 \\
	\hline

	8                                                                                  &
	factorial\_tail(1, 2, Res). \newline
	Успешная унификация. \newline
	Вычисление предиката is.                                                           &
	NewN is 2 - 1 = \newline
	Истина \newline
	Унификация успешна. \newline
	Подстановка: \newline
	\{NewN = 1\}                                                                       &
	Унификация успешна. \newline
	Резольвента непустая. \newline
	Прямой ход. \newline
	Запуск алгоритма унификации для factorial\_tail(1, 2, Res) и 1-ого заголовка БЗ.           \\
	\hline

	9                                                                                  &
	factorial\_tail(1, 2, Res).                                                        &
	factorial\_tail(1, 2, Res) = \newline
	factorial\_tail(0, Acc, Acc) \newline
	Унификация неуспешна (разные константы: 1 и 0).                                    &
	Унификация неуспешна. \newline
	Резольвента непустая. \newline
	Прямой ход. Переход к следующему правилу.                                                  \\
	\hline

	10                                                                                 &
	>(1, 0), \newline
	NewAcc is 2 * 1, \newline
	NewN is 1 - 1, \newline
	factorial\_tail(NewN, NewAcc, Res). \newline
	Успешная унификация. \newline
	Редукция и применение подстановки.                                                 &
	factorial\_tail(1, 2, Res) = \newline
	factorial\_tail(N, Acc, Result) \newline
	Унификация успешна. \newline
	Подстановка: \newline
	\{N = 1, Acc = 2, Result = Res\}                                                   &
	Унификация успешна. \newline
	Резольвента непустая. \newline
	Прямой ход. \newline
	Запуск алгоритма унификации для >(1, 0) и Истина.                                          \\
	\hline

	11                                                                                 &
	NewAcc is 2 * 1, \newline
	NewN is 1 - 1, \newline
	factorial\_tail(NewN, NewAcc, Res). \newline
	Успешная унификация.                                                               &
	>(1, 0) = \newline
	Истина \newline
	Унификация успешна. \newline
	Подстановка: \newline
	Нет.                                                                               &
	Унификация успешна. \newline
	Резольвента непустая. \newline
	Прямой ход.                                                                                \\
	\hline

	12                                                                                 &
	NewN is 1 - 1, \newline
	factorial\_tail(NewN, 2, Res). \newline
	Успешная унификация.                                                               &
	NewAcc is 2 * 1 = \newline
	Истина \newline
	Унификация успешна. \newline
	Подстановка: \newline
	\{NewAcc = 2\}                                                                     &
	Унификация успешна. \newline
	Резольвента непустая. \newline
	Прямой ход.                                                                                \\
	\hline

	13                                                                                 &
	factorial\_tail(0, 2, Res). \newline
	Успешная унификация.                                                               &
	NewN is 1 - 1 = \newline
	Истина \newline
	Унификация успешна. \newline
	Подстановка: \newline
	\{NewN = 0\}                                                                       &
	Унификация успешна. \newline
	Резольвента непустая. \newline
	Прямой ход. \newline
	Запуск алгоритма унификации для factorial\_tail(0, 2, Res) и 1-ого заголовка БЗ.           \\
	\hline

	14                                                                                 &
	!. \newline
	Успешная унификация. \newline
	Редукция базового случая (дно рекурсии). \newline
	Применение подстановки.                                                            &
	factorial\_tail(0, 2, Res) = \newline
	factorial\_tail(0, Acc, Acc) \newline
	Унификация успешна. \newline
	Подстановка: \newline
	\{Acc = 2, Res = 2\}                                                               &
	Унификация успешна. \newline
	Резольвента непустая. \newline
	Прямой ход. \newline
	Запуск алгоритма обработки предиката отсечения (!).                                        \\
	\hline

	15                                                                                 &
	Резольвента пуста. \newline
	Успешная унификация. \newline
	Отработка отсечения.                                                               &
	! = \newline
	! \newline
	Унификация успешна. \newline
	Подстановка: \newline
	Нет.                                                                               &
	Унификация успешна. \newline
	Предикат отсечения срабатывает: удаляются точки возврата. \newline
	Резольвента пуста. \newline
	Знания подобраны. Подстановка \{Res = 2\}. \newline
	Завершение работы.                                                                         \\
	\hline
\end{tabularx}

\clearpage
Вопрос: \texttt{?- fibonacci(2, Res).}

\begin{tabularx}{\textwidth}{|c|>{\small}X|>{\small}X|>{\small}X|}
	\hline
	\textbf{№}                                                                         &
	\textbf{Состояние резольвенты}                                                     &
	\textbf{Для каких термов запускается алгоритм унификации, результат и подстановка} &
	\textbf{Дальнейшие действия}                                                         \\
	\hline

	0                                                                                  &
	fibonacci(2, Res).                                                                 &
	—                                                                                  &
	Резольвента непустая. \newline
	Прямой ход. \newline
	Запуск алгоритма унификации для верхней подцели с 1-ым заголовком из БЗ.             \\
	\hline

	1                                                                                  &
	fibonacci(2, Res).                                                                 &
	fibonacci(2, Res) = \newline
	factorial\_tail(0, Acc, Acc) \newline
	Унификация неуспешна (разные функторы).                                            &
	Унификация неуспешна. \newline
	Резольвента непустая. \newline
	Прямой ход. Переход к следующему правилу.                                            \\
	\hline

	2                                                                                  &
	...                                                                                &
	... (Пропуск остальных предикатов factorial\_tail, factorial, fib\_tail)           &
	...                                                                                  \\
	\hline

	3                                                                                  &
	fib\_tail(2, 0, 1, Res). \newline
	Успешная унификация. \newline
	Редукция и применение подстановки к полученной конъюнкции целей: \newline
	fib\_tail(2, 0, 1, Res).                                                           &
	fibonacci(2, Res) = \newline
	fibonacci(N, Result) \newline
	Унификация успешна. \newline
	Подстановка: \newline
	\{N = 2, Result = Res\}                                                            &
	Унификация успешна. \newline
	Резольвента непустая. \newline
	Прямой ход. \newline
	Запуск алгоритма унификации для fib\_tail(2, 0, 1, Res) и 1-ого заголовка БЗ.        \\
	\hline

	4                                                                                  &
	fib\_tail(2, 0, 1, Res).                                                           &
	fib\_tail(2, 0, 1, Res) = \newline
	fib\_tail(0, Acc1, \_, Acc1) \newline
	Унификация неуспешна (разные константы: 2 и 0).                                    &
	Унификация неуспешна. \newline
	Резольвента непустая. \newline
	Прямой ход. Переход к следующему правилу.                                            \\
	\hline

	5                                                                                  &
	>(2, 0), \newline
	NewAcc2 is 0 + 1, \newline
	NewN is 2 - 1, \newline
	fib\_tail(NewN, 1, NewAcc2, Res). \newline
	Успешная унификация. \newline
	Редукция и применение подстановки.                                                 &
	fib\_tail(2, 0, 1, Res) = \newline
	fib\_tail(N, Acc1, Acc2, Result) \newline
	Унификация успешна. \newline
	Подстановка: \newline
	\{N = 2, Acc1 = 0, Acc2 = 1, Result = Res\}                                        &
	Унификация успешна. \newline
	Резольвента непустая. \newline
	Прямой ход. \newline
	Запуск алгоритма унификации для >(2, 0) и Истина.                                    \\
	\hline

	6                                                                                  &
	NewAcc2 is 0 + 1, \newline
	NewN is 2 - 1, \newline
	fib\_tail(NewN, 1, NewAcc2, Res). \newline
	Успешная унификация. \newline
	Вычисление предиката сравнения.                                                    &
	>(2, 0) = \newline
	Истина \newline
	Унификация успешна. \newline
	Подстановка: \newline
	Нет.                                                                               &
	Унификация успешна. \newline
	Резольвента непустая. \newline
	Прямой ход. \newline
	Запуск алгоритма вычисления для is(0 + 1).                                           \\
	\hline

	7                                                                                  &
	NewN is 2 - 1, \newline
	fib\_tail(NewN, 1, 1, Res). \newline
	Успешная унификация. \newline
	Вычисление предиката is.                                                           &
	NewAcc2 is 0 + 1 = \newline
	Истина \newline
	Унификация успешна. \newline
	Подстановка: \newline
	\{NewAcc2 = 1\}                                                                    &
	Унификация успешна. \newline
	Резольвента непустая. \newline
	Прямой ход. \newline
	Запуск алгоритма вычисления для is(2 - 1).                                           \\
	\hline

	8                                                                                  &
	fib\_tail(1, 1, 1, Res). \newline
	Успешная унификация. \newline
	Вычисление предиката is.                                                           &
	NewN is 2 - 1 = \newline
	Истина \newline
	Унификация успешна. \newline
	Подстановка: \newline
	\{NewN = 1\}                                                                       &
	Унификация успешна. \newline
	Резольвента непустая. \newline
	Прямой ход. \newline
	Запуск алгоритма унификации для fib\_tail(1, 1, 1, Res) и 1-ого заголовка БЗ.        \\
	\hline

	9                                                                                  &
	fib\_tail(1, 1, 1, Res).                                                           &
	fib\_tail(1, 1, 1, Res) = \newline
	fib\_tail(0, Acc1, \_, Acc1) \newline
	Унификация неуспешна (разные константы: 1 и 0).                                    &
	Унификация неуспешна. \newline
	Резольвента непустая. \newline
	Прямой ход. Переход к следующему правилу.                                            \\
	\hline

	10                                                                                 &
	>(1, 0), \newline
	NewAcc2 is 1 + 1, \newline
	NewN is 1 - 1, \newline
	fib\_tail(NewN, 1, NewAcc2, Res). \newline
	Успешная унификация. \newline
	Редукция и применение подстановки.                                                 &
	fib\_tail(1, 1, 1, Res) = \newline
	fib\_tail(N, Acc1, Acc2, Result) \newline
	Унификация успешна. \newline
	Подстановка: \newline
	\{N = 1, Acc1 = 1, Acc2 = 1, Result = Res\}                                        &
	Унификация успешна. \newline
	Резольвента непустая. \newline
	Прямой ход. \newline
	Запуск алгоритма унификации для >(1, 0) и Истина.                                    \\
	\hline

	11                                                                                 &
	NewAcc2 is 1 + 1, \newline
	NewN is 1 - 1, \newline
	fib\_tail(NewN, 1, NewAcc2, Res). \newline
	Успешная унификация.                                                               &
	>(1, 0) = \newline
	Истина \newline
	Унификация успешна. \newline
	Подстановка: \newline
	Нет.                                                                               &
	Унификация успешна. \newline
	Резольвента непустая. \newline
	Прямой ход.                                                                          \\
	\hline

	12                                                                                 &
	NewN is 1 - 1, \newline
	fib\_tail(NewN, 1, 2, Res). \newline
	Успешная унификация.                                                               &
	NewAcc2 is 1 + 1 = \newline
	Истина \newline
	Унификация успешна. \newline
	Подстановка: \newline
	\{NewAcc2 = 2\}                                                                    &
	Унификация успешна. \newline
	Резольвента непустая. \newline
	Прямой ход.                                                                          \\
	\hline

	13                                                                                 &
	fib\_tail(0, 1, 2, Res). \newline
	Успешная унификация.                                                               &
	NewN is 1 - 1 = \newline
	Истина \newline
	Унификация успешна. \newline
	Подстановка: \newline
	\{NewN = 0\}                                                                       &
	Унификация успешна. \newline
	Резольвента непустая. \newline
	Прямой ход. \newline
	Запуск алгоритма унификации для fib\_tail(0, 1, 2, Res) и 1-ого заголовка БЗ.        \\
	\hline

	14                                                                                 &
	!. \newline
	Успешная унификация. \newline
	Редукция базового случая (дно рекурсии). \newline
	Применение подстановки.                                                            &
	fib\_tail(0, 1, 2, Res) = \newline
	fib\_tail(0, Acc1, \_, Acc1) \newline
	Унификация успешна. \newline
	Подстановка: \newline
	\{Acc1 = 1, \_ = 2, Res = 1\}                                                      &
	Унификация успешна. \newline
	Резольвента непустая. \newline
	Прямой ход. \newline
	Запуск алгоритма обработки предиката отсечения (!).                                  \\
	\hline

	15                                                                                 &
	Резольвента пуста. \newline
	Успешная унификация. \newline
	Отработка отсечения.                                                               &
	! = \newline
	! \newline
	Унификация успешна. \newline
	Подстановка: \newline
	Нет.                                                                               &
	Унификация успешна. \newline
	Предикат отсечения срабатывает: удаляются точки возврата. \newline
	Резольвента пуста. \newline
	Знания подобраны. Подстановка \{Res = 1\}. \newline
	Завершение работы.                                                                   \\
	\hline
\end{tabularx}
